\documentclass{article}

\usepackage[accepted]{icml2025}
\usepackage{microtype}
\usepackage{graphicx}
\usepackage{subcaption}
\usepackage{booktabs}
\usepackage{tabularx}
\usepackage{array}
\usepackage{amsmath}
\usepackage{amssymb}
\usepackage{algorithm}
\usepackage{algorithmic}
\usepackage{url}
\usepackage{hyperref}

\newcommand{\safeincludegraphics}[2][]{%
  \IfFileExists{#2}{%
    \includegraphics[#1]{#2}%
  }{%
    \fbox{\parbox[c][1.35in][c]{0.88\linewidth}{\centering Missing figure: \texttt{#2}}}%
  }%
}
\newcolumntype{Y}{>{\raggedright\arraybackslash}X}

\icmltitlerunning{A Small-Scale Study of Muon and Newton--Schulz Iterations}

\begin{document}

\twocolumn[
\icmltitle{A Small-Scale Experimental Study of the Muon Optimizer and Newton--Schulz Iterations}

\begin{icmlauthorlist}
\icmlauthor{Zheng Renjie}{sustech}
\icmlauthor{Chen Ximing}{sustech}
\icmlauthor{Jiang Hongjun}{sustech}
\end{icmlauthorlist}

\icmlaffiliation{sustech}{Southern University of Science and Technology, Shenzhen, China}
\icmlcorrespondingauthor{Zheng Renjie}{12412840@mail.sustech.edu.cn}

\icmlkeywords{Optimization, Muon, AdamW, Newton--Schulz, Polar Decomposition, Image Classification}

\vskip 0.3in
]

\begin{abstract}
This project studies whether the Muon optimizer can provide faster and more stable training than AdamW on small neural networks, and how its behavior depends on the Newton--Schulz approximation used to compute matrix-orthogonalized updates. We implement a PyTorch training framework with AdamW and Muon, evaluate MLP and small CNN models on MNIST and Fashion-MNIST, run equal-budget hyperparameter search, and ablate the Newton--Schulz iteration count and coefficients. On Fashion-MNIST with an MLP, Muon improves final validation accuracy from $88.75\%\pm0.49\%$ to $89.95\%\pm0.28\%$ over five seeds, while reducing final training loss from $0.2063$ to $0.0306$. With a small CNN, Muon improves accuracy from $86.57\%\pm0.71\%$ to $90.78\%\pm0.65\%$. Newton--Schulz ablations show that the practical Muon coefficients converge faster than a standard polynomial choice, but large iteration counts can over-optimize the training loss without improving final validation accuracy. These results suggest that Muon is a fast and useful optimizer for matrix-valued hidden layers, but its benefits depend strongly on the iteration design and early-stopping behavior.
\end{abstract}

\section{Introduction}

Optimization is a central problem in modern deep learning. Adam and AdamW are widely used because adaptive coordinate-wise scaling gives stable performance across many tasks \citep{kingma2015adam,loshchilov2019decoupled}. Muon is a more recent optimizer designed for hidden layers of neural networks \citep{jordan2024muon}. Instead of applying a coordinate-wise adaptive update, Muon treats hidden-layer weights as matrices, applies momentum, and then approximately orthogonalizes the momentum matrix before updating the parameters.

The mathematical core of Muon is connected to polar decomposition and Newton--Schulz iterations \citep{higham2008functions,kim2026convergence}. This makes it a useful topic for a mathematical foundations project: the optimizer is not only a practical training method, but also an example of how numerical linear algebra changes optimization dynamics.

This report asks three questions:
\begin{enumerate}
    \item Can Muon converge faster or achieve better validation accuracy than AdamW on small neural networks?
    \item How sensitive is Muon to the Newton--Schulz iteration count $q$?
    \item Do the practical Muon coefficients behave differently from a standard Newton--Schulz polynomial choice?
\end{enumerate}

The experiments intentionally use small and reproducible image-classification tasks rather than large language model training. This keeps the project computationally controlled while still exposing the main optimizer behavior.

\subsection{Problem Formulation}

All experiments solve empirical risk minimization for supervised classification:
\begin{equation}
\min_{\theta}\; \mathcal{L}(\theta)
= \frac{1}{n}\sum_{i=1}^{n}\ell(f_\theta(x_i),y_i),
\label{eq:erm}
\end{equation}
where $f_\theta$ is an MLP or CNN, $\ell$ is cross-entropy loss, and $\theta$ contains matrix-valued hidden parameters as well as non-matrix parameters such as classifier-head weights. The project compares two update rules for minimizing Eq.~\eqref{eq:erm}: AdamW, which rescales coordinates using first- and second-moment estimates, and Muon, which applies a matrix-valued polar-direction update to hidden layers. The key experimental question is whether this matrix update improves optimization speed, validation accuracy, or stability under the same training budget.

\section{Background}

\subsection{AdamW}

AdamW maintains exponential moving averages of the stochastic gradient and squared gradient. Given gradient $g_t$, it updates
\begin{align}
m_t &= \beta_1 m_{t-1} + (1-\beta_1)g_t,\\
v_t &= \beta_2 v_{t-1} + (1-\beta_2)g_t \odot g_t.
\end{align}
After bias correction, AdamW applies a coordinate-wise update and decoupled weight decay \citep{loshchilov2019decoupled}. This makes it a strong baseline because it is robust, easy to tune, and commonly used in neural network training.

\subsection{Muon}

Muon uses a matrix-geometric update for hidden-layer weights. For a matrix parameter $W_t$ with gradient $G_t$, Muon first computes a momentum matrix
\begin{equation}
M_t = \beta M_{t-1} + (1-\beta)G_t.
\end{equation}
Ideally, the update direction is the polar factor
\begin{equation}
\operatorname{Polar}(M_t) = U_t V_t^\top,
\end{equation}
where $M_t = U_t \Sigma_t V_t^\top$ is an SVD. Computing an SVD at every optimizer step would be expensive, so practical Muon approximates this polar direction with a few Newton--Schulz iterations.

In this project, the Newton--Schulz update has the form
\begin{equation}
X_{j+1} = aX_j + \left(bX_jX_j^\top + c(X_jX_j^\top)^2\right)X_j,
\label{eq:ns}
\end{equation}
for $j=0,\ldots,q-1$. We compare the practical Muon coefficients
\[
(a,b,c)=(3.4445,-4.7750,2.0315)
\]
with a more standard polynomial choice
\[
(a,b,c)=\left(\frac{15}{8},-\frac{5}{4},\frac{3}{8}\right).
\]

\begin{algorithm}[t]
\caption{Muon optimizer step}
\label{alg:muon}
\begin{algorithmic}[1]
\REQUIRE Hidden matrix parameter $W$, gradient $G$, momentum $M$, learning rate $\alpha$, momentum coefficient $\beta$, Newton--Schulz steps $q$, coefficients $(a,b,c)$
\STATE $M \leftarrow \beta M + (1-\beta)G$
\STATE Reshape $M$ to a matrix $\tilde M\in\mathbb{R}^{d_1\times d_2}$ if $W$ is convolutional
\STATE Normalize $X_0 \leftarrow \tilde M / (\|\tilde M\|_F+\epsilon)$
\FOR{$j=0$ to $q-1$}
\STATE $X_{j+1} \leftarrow aX_j + (bX_jX_j^\top+c(X_jX_j^\top)^2)X_j$
\ENDFOR
\STATE Reshape $X_q$ back to the shape of $W$
\STATE $W \leftarrow (1-\alpha\lambda)W - \alpha X_q$
\STATE Update non-matrix parameters separately with AdamW
\end{algorithmic}
\end{algorithm}

\section{Experimental Setup}

\subsection{Tasks and Models}

We evaluate supervised image classification on MNIST \citep{lecun1998mnist} and Fashion-MNIST \citep{xiao2017fashion}. The main controlled setting is an MLP with two hidden linear layers. We also run a small CNN robustness check on Fashion-MNIST.

The MLP uses bias-free matrix layers so that optimizer behavior is easier to interpret. The small CNN uses two convolutional blocks followed by an adaptive pooling layer and a linear classifier. All models use cross-entropy loss.

\subsection{Optimizer Protocol}

AdamW is applied to all trainable parameters. Muon is applied only to matrix-valued hidden parameters: hidden linear weights and reshaped convolutional weights. Biases, normalization-like parameters, and classifier-head parameters are excluded from Muon and handled by AdamW when present. This follows the intended use of Muon as an optimizer for hidden layers rather than scalar or head parameters.

For the main MLP experiments, AdamW uses learning rate $10^{-3}$ and weight decay $10^{-4}$. Muon uses learning rate $10^{-2}$, weight decay $10^{-4}$, $q=3$, and the practical Muon coefficients. We also run an equal-budget search on Fashion-MNIST with 12 AdamW configurations and 12 Muon configurations. The search supports these representative choices: the best AdamW configuration is $(\alpha,\lambda)=(3\cdot10^{-3},10^{-4})$, while the best Muon configuration is $(10^{-2},10^{-4})$.

\subsection{Metrics}

We record training loss, validation loss, validation accuracy, epoch time, and peak GPU memory. Main comparisons are averaged over multiple random seeds. MNIST uses three seeds. Fashion-MNIST MLP and small CNN experiments use five seeds. Hyperparameter search uses one seed and is used only to verify the relative tuning budget, not as a final statistical estimate.

\subsection{Reproducibility Details}

The implementation is written in PyTorch \citep{paszke2019pytorch}. Every run writes a CSV row for each epoch with dataset, model, optimizer, learning rate, weight decay, seed, training loss, validation loss, validation accuracy, epoch time, peak GPU memory, Newton--Schulz steps, and coefficient set. All figures and tables are generated from these logs. Before the full experiments, I verified that AdamW and Muon each complete a one-epoch smoke test, that fixed seeds reproduce the same setup, and that Muon receives only matrix-valued hidden parameters. Experiments were run on an NVIDIA RTX PRO 6000 Blackwell Server Edition GPU with batch size 512. MNIST uses 10 epochs, Fashion-MNIST MLP uses 15 epochs, the small CNN main run uses 12 epochs, and small CNN ablations use 10 epochs.

\section{Results}

\subsection{MNIST Sanity Check}

MNIST is an easy benchmark, so it is mainly used to verify that both optimizers run correctly. Table~\ref{tab:mnist} shows that AdamW and Muon both reach high validation accuracy. Muon has slightly higher final accuracy and much lower final training loss, but it is slower per epoch.

\begin{table}[t]
\centering
\caption{MNIST sanity check.}
\label{tab:mnist}
\small
\setlength{\tabcolsep}{3pt}
\resizebox{\columnwidth}{!}{%
\begin{tabular}{lrrrr}
\toprule
Method & Acc. & Val Loss & Train Loss & Time \\
\midrule
AdamW & $98.05\pm0.17$ & 0.0706 & 0.0159 & 2.566 \\
Muon $q=3$ & $98.39\pm0.09$ & 0.0728 & 0.0025 & 3.264 \\
\bottomrule
\end{tabular}
}
\end{table}

\subsection{Fashion-MNIST MLP}

Fashion-MNIST is the main experiment because it is harder than MNIST but still small enough for repeated runs. Table~\ref{tab:main} and Figure~\ref{fig:main_acc} show that Muon improves final validation accuracy from $88.75\%$ to $89.95\%$ over five seeds. The final training loss is also much lower for Muon, decreasing from $0.2063$ to $0.0306$.

\begin{table}[t]
\centering
\caption{Fashion-MNIST MLP main comparison.}
\label{tab:main}
\small
\setlength{\tabcolsep}{3pt}
\resizebox{\columnwidth}{!}{%
\begin{tabular}{lrrrr}
\toprule
Method & Acc. & Val Loss & Train Loss & Time \\
\midrule
AdamW & $88.75\pm0.49$ & 0.3329 & 0.2063 & 2.558 \\
Muon $q=3$ & $89.95\pm0.28$ & 0.4424 & 0.0306 & 2.784 \\
\bottomrule
\end{tabular}
}
\end{table}

\begin{figure}[t]
\centering
\safeincludegraphics[width=\columnwidth]{figures/main_final_accuracy_bar.png}
\caption{Fashion-MNIST MLP final accuracy.}
\label{fig:main_acc}
\end{figure}

The validation loss tells a more nuanced story. Table~\ref{tab:best_epoch} separates final-epoch metrics from best-epoch behavior. Muon reaches its best validation loss around epoch 5.2, while AdamW reaches it around epoch 13.0. Thus, Muon optimizes faster, but continuing training can overfit. This is also visible in the time-to-threshold metric: Muon reaches the selected training-loss threshold around epochs 8--9, while AdamW does not reach it within 15 epochs.

\begin{table}[t]
\centering
\caption{Fashion-MNIST MLP best-epoch metrics.}
\label{tab:best_epoch}
\small
\setlength{\tabcolsep}{4pt}
\resizebox{\columnwidth}{!}{%
\begin{tabular}{lrr}
\toprule
Metric & AdamW & Muon $q=3$ \\
\midrule
Best val. acc. & 89.09 & 90.25 \\
Best acc. epoch & 13.4 & 12.8 \\
Best val. loss & 0.3187 & 0.2970 \\
Best loss epoch & 13.0 & 5.2 \\
\bottomrule
\end{tabular}
}
\end{table}

\subsection{Hyperparameter Search}

Table~\ref{tab:search} reports the top configurations from the equal-budget search. The search is not intended as a final multi-seed comparison, but it reduces the risk that the main result comes from an unfair choice of learning rates. The best AdamW setting reaches $88.02\%$ after eight epochs, while the best Muon setting reaches $90.25\%$.

\begin{table}[t]
\centering
\caption{Fashion-MNIST hyperparameter search.}
\label{tab:search}
\resizebox{\columnwidth}{!}{%
\begin{tabular}{lrrrr}
\toprule
Method & LR & WD & Acc. & Val Loss \\
\midrule
AdamW & 0.003 & 0.0001 & 88.02 & 0.3452 \\
AdamW & 0.003 & 0.0100 & 87.69 & 0.3513 \\
AdamW & 0.003 & 0.0000 & 87.59 & 0.3583 \\
Muon $q=3$ & 0.010 & 0.0001 & 90.25 & 0.3033 \\
Muon $q=3$ & 0.010 & 0.0100 & 90.24 & 0.2997 \\
Muon $q=3$ & 0.010 & 0.0000 & 90.14 & 0.3045 \\
\bottomrule
\end{tabular}}
\end{table}

\begin{figure}[t]
\centering
\safeincludegraphics[width=\columnwidth]{figures/search_top_configs_bar.png}
\caption{Hyperparameter search top configurations.}
\label{fig:search}
\end{figure}

\subsection{Newton--Schulz Ablation on MLP}

The Newton--Schulz ablation is the main mathematical experiment. Table~\ref{tab:ablation_mlp} compares $q\in\{1,2,3,5\}$ and two coefficient choices. With the practical Muon coefficients, increasing $q$ sharply reduces training loss: from $0.2150$ at $q=1$ to $0.0062$ at $q=5$. However, final validation accuracy does not improve beyond $q=3$, and final validation loss becomes much worse for $q=5$.

\begin{table*}[t]
\centering
\caption{Fashion-MNIST MLP Newton--Schulz ablation.}
\label{tab:ablation_mlp}
\begin{tabular}{llrrrr}
\toprule
Coeff. & $q$ & Acc. & Val Loss & Train Loss & Time \\
\midrule
Muon & 1 & $88.84\pm0.13$ & 0.3213 & 0.2150 & 2.779 \\
Muon & 2 & $89.67\pm0.21$ & 0.3392 & 0.1093 & 2.817 \\
Muon & 3 & $89.95\pm0.28$ & 0.4424 & 0.0306 & 2.843 \\
Muon & 5 & $89.85\pm0.10$ & 0.5652 & 0.0062 & 3.162 \\
Standard & 1 & $88.15\pm0.20$ & 0.3317 & 0.2556 & 2.882 \\
Standard & 2 & $88.86\pm0.21$ & 0.3203 & 0.2118 & 2.739 \\
Standard & 3 & $89.33\pm0.16$ & 0.3196 & 0.1625 & 2.888 \\
Standard & 5 & $89.80\pm0.20$ & 0.3814 & 0.0661 & 3.047 \\
\bottomrule
\end{tabular}
\end{table*}

\begin{figure}[t]
\centering
\safeincludegraphics[width=\columnwidth]{figures/ablation_final_accuracy_bar.png}
\caption{MLP ablation final accuracy.}
\label{fig:ablation_acc}
\end{figure}

\begin{figure}[t]
\centering
\safeincludegraphics[width=\columnwidth]{figures/ablation_loss_tradeoff.png}
\caption{MLP ablation loss trade-off.}
\label{fig:ablation_loss}
\end{figure}

The standard coefficients are more conservative. Their training loss decreases more slowly, and validation loss remains more stable. This supports the interpretation that the practical coefficients are not merely approximating the polar factor more accurately; they change the optimization dynamics by producing stronger updates.

\subsection{Small CNN Robustness Check}

To check whether the MLP result is architecture-specific, we repeat the AdamW vs Muon comparison with a small CNN. Table~\ref{tab:cnn_main} and Figure~\ref{fig:cnn_main} show a stronger result: Muon improves validation accuracy from $86.57\%$ to $90.78\%$. In this setting, Muon also improves validation loss, unlike the final-epoch MLP result.

\begin{table}[t]
\centering
\caption{Small CNN main comparison.}
\label{tab:cnn_main}
\small
\setlength{\tabcolsep}{3pt}
\resizebox{\columnwidth}{!}{%
\begin{tabular}{lrrrr}
\toprule
Method & Acc. & Val Loss & Train Loss & Time \\
\midrule
AdamW & $86.57\pm0.71$ & 0.3821 & 0.3497 & 2.953 \\
Muon $q=3$ & $90.78\pm0.65$ & 0.2677 & 0.1831 & 3.240 \\
\bottomrule
\end{tabular}
}
\end{table}

\begin{figure}[t]
\centering
\safeincludegraphics[width=\columnwidth]{figures/small_cnn_final_accuracy_bar.png}
\caption{Small CNN final accuracy.}
\label{fig:cnn_main}
\end{figure}

The CNN ablation in Table~\ref{tab:cnn_ablation} gives a related but not identical picture to the MLP. Practical Muon coefficients still work best around $q=3$, while standard coefficients improve steadily with larger $q$. The result suggests that the useful iteration count depends on architecture and optimization regime.

\begin{table*}[t]
\centering
\caption{Small CNN ablation table.}
\label{tab:cnn_ablation}
\begin{tabular}{llrrrr}
\toprule
Coeff. & $q$ & Acc. & Val Loss & Train Loss & Time \\
\midrule
Muon & 1 & $88.87\pm0.49$ & 0.3138 & 0.2796 & 3.342 \\
Muon & 2 & $89.95\pm0.35$ & 0.2873 & 0.2312 & 3.254 \\
Muon & 3 & $90.61\pm0.23$ & 0.2683 & 0.2004 & 3.325 \\
Muon & 5 & $90.40\pm0.42$ & 0.2751 & 0.1963 & 3.385 \\
Standard & 1 & $87.89\pm0.78$ & 0.3435 & 0.3121 & 3.257 \\
Standard & 2 & $88.81\pm0.47$ & 0.3168 & 0.2825 & 3.436 \\
Standard & 3 & $89.56\pm0.53$ & 0.2986 & 0.2542 & 3.479 \\
Standard & 5 & $90.34\pm0.56$ & 0.2776 & 0.2160 & 3.458 \\
\bottomrule
\end{tabular}
\end{table*}

\begin{figure}[t]
\centering
\safeincludegraphics[width=\columnwidth]{figures/small_cnn_ablation_accuracy_bar.png}
\caption{Small CNN ablation accuracy.}
\label{fig:cnn_ablation}
\end{figure}

\section{Discussion}

The experiments support four findings.

First, Muon consistently lowers training loss faster than AdamW. This is most visible on Fashion-MNIST MLP, where Muon reaches the selected training-loss threshold around epochs 8--9, while AdamW does not reach it within 15 epochs.

Second, faster training loss reduction does not always imply lower final validation loss. On the MLP, $q=5$ with practical Muon coefficients drives training loss close to zero but worsens validation loss. This negative result is important: it shows that a more aggressive or more repeated Newton--Schulz update is not automatically better. Muon can benefit from early stopping, adjusted regularization, or a smaller iteration count.

Third, Newton--Schulz coefficients matter. Practical Muon coefficients reach high validation accuracy with fewer iterations, while standard coefficients are more conservative and often need larger $q$. Therefore, the coefficient choice is part of the optimizer design, not only a numerical approximation detail.

Fourth, the small CNN experiment shows that the result is not limited to MLPs. In fact, Muon gives a larger gain on the CNN in these runs. This may be because convolutional kernels are natural matrix-valued hidden parameters after reshaping, making them suitable for Muon's matrix update. The CNN ablation is also a mixed result: unlike the MLP, increasing $q$ to 5 remains competitive for the standard coefficients, suggesting that the best Newton--Schulz design is architecture-dependent.

\section{Limitations}

This project is intentionally small-scale. The results do not prove that Muon is always better than AdamW, nor do they directly transfer to large language model training. The hyperparameter search uses one seed, so it should be treated as a fairness check rather than a final benchmark. The experiments also use fixed training budgets; some Muon configurations, especially large $q$, may look better under early stopping. Finally, the implementation uses a practical subset of Muon behavior and does not analyze all possible parameter-group protocols.

\section{Conclusion}

This project implemented a reproducible PyTorch study of AdamW, Muon, and Newton--Schulz iterations on small neural networks. Muon achieved faster loss reduction and higher validation accuracy than AdamW on Fashion-MNIST MLP and small CNN models. The Newton--Schulz ablations show that $q$ and coefficient choice strongly affect the trade-off among convergence speed, validation accuracy, and computational cost. In the tested settings, $q=3$ with practical Muon coefficients is a strong default: it is faster and more accurate than small $q$, while avoiding some of the overfitting behavior seen at $q=5$. Overall, Muon is a promising optimizer for matrix-valued hidden layers, but its practical success depends on careful iteration design and evaluation beyond final training loss.

\bibliography{example_paper}
\bibliographystyle{icml2025}

\end{document}
